\section*{Chapter \ref{ch:b2:field-lab-methods-2} --- Field and Lab Methods II}
% numbers from the Book 4 second-writer scratch calculator calc/ch31.py

\begin{solution}{exo:b2:field-lab-methods-2:1}
(a) (i) \omterm{def:b2:field-lab-methods-2:textures}{Ophitic texture}; (ii) \omterm{def:b2:field-lab-methods-2:textures}{granoblastic texture}; (iii) a \omterm{def:b2:field-lab-methods-2:porphyroblast}{poikiloblast}, a \omterm{def:b2:field-lab-methods-2:porphyroblast}{porphyroblast} full of
inclusions; (iv) a \omterm{def:b2:field-lab-methods-2:porphyroblast}{corona}, a reaction rim around a garnet that was breaking down. (b) By
\eqref{eq:b2:field-lab-methods-2:lami}, $\gamma_1/\gamma_2 = \sin\ang{104}/\sin\ang{128} =
0.9703/0.7880 = 1.23$: that boundary pulls \qty{23}{\percent} harder than the two others.
\end{solution}

\begin{solution}{exo:b2:field-lab-methods-2:2}
(a) $Q + A + P = 80$: Q $= 100 \times 25/80 = 31.3$, A $= 18.8$, P $= 50.0$. (b) P/(A+P) $=
40/55 = \qty{72.7}{\percent}$, between 65 and 90 with Q between 20 and 60: a granodiorite.
(c) $M' = 12 + 8 = \qty{20}{\percent}$, below 35: leucocratic.
\end{solution}

\begin{solution}{exo:b2:field-lab-methods-2:3}
(a) $\sigma = \sqrt{0.12 \times 0.88/800} = 0.0115$, that is $12.0 \pm 1.1$ percentage points, a
relative error of $0.0115/0.12 = \qty{9.6}{\percent}$. (b) Four times the points halve it: $0.0057$.
\end{solution}

\begin{solution}{exo:b2:field-lab-methods-2:4}
(a) The loss is \qty{0.0288}{g}, and $0.0288/0.8000 = \qty{3.60}{\percent}$.
(b) $0.7712/8.7712 = 0.0879$.
(c) Ignited powder: $5.20/0.08792 = \qty{59.1}{\percent}$; original powder:
$59.14 \times (1 - 0.0360) = \qty{57.0}{\percent}$.
\end{solution}

\begin{solution}{exo:b2:field-lab-methods-2:5}
(a) $\sigma = 0.10 \times 0.02 = 0.002$ and $n = 0.02 \times 0.98/0.002^2 = 4900$ points.
(b) $\sigma = 0.03$ and $n = 0.30 \times 0.70/0.03^2 = 233.3$, so 234 points. (c) For a fixed relative
error, $n = (1-p)/(p\,\varepsilon^2)$ grows as $1/p$: the accessory mineral demands 21 times more
points, more than a section usually provides.
\end{solution}

\begin{solution}{exo:b2:field-lab-methods-2:6}
(a) RGM-1: $72.89/73.45 = \qty{99.2}{\percent}$; BCR-2: $54.17/54.1 = \qty{100.1}{\percent}$.
(b) RGM-1: $0.19/\sqrt{20} = 0.042$, and $-0.56/0.042 = -13.2$; BCR-2: $0.13/\sqrt{20} = 0.029$, and
$0.07/0.029 = 2.4$. (c) RGM-1 is biased: its mean misses the recommended value by 13 times its own
standard deviation, far more than chance allows (if the recommended value is right). BCR-2 is at
the limit of what chance gives. RGM-1 is not less precise: its relative standard deviation,
$0.19/72.89 = \qty{0.26}{\percent}$, is close to that of BCR-2, \qty{0.24}{\percent}. The problem is
trueness, not precision.
\end{solution}

\begin{solution}{exo:b2:field-lab-methods-2:7}
(a) $M = 39.098 + 3 \times 24.305 + 26.982 + 3 \times 28.085 + 12 \times 15.999 + 2 \times 1.008 =
\qty{417.25}{g.mol^{-1}}$; the two OH make one \ce{H2O} ($18.015$), \qty{4.32}{\percent}.
(b) \ce{K2O} $47.10/417.25 = \qty{11.3}{\percent}$; \ce{MgO} $3 \times 40.304/417.25 =
\qty{29.0}{\percent}$; \ce{Al2O3} $50.98/417.25 = \qty{12.2}{\percent}$; \ce{SiO2} $3 \times
60.083/417.25 = \qty{43.2}{\percent}$; total \qty{95.7}{\percent}. (c) The microprobe does not
detect hydrogen: a mica analysis lacks its water, about \qty{4}{\percent}, and should total near
96.
\end{solution}

\begin{solution}{exo:b2:field-lab-methods-2:8}
(a) $\beta = \ln(0.1194/0.11905)/\ln(85.909261/87.905612) = 0.002936/(-0.02297) = -0.1278$.
(b) $0.71062 \times (86.908878/85.909261)^{-0.1278} = 0.70957$. (c) The measured 86/88 is below the
true ratio: the heavy isotope was favoured, and the measured 87/86 too high.
\end{solution}

\begin{solution}{exo:b2:field-lab-methods-2:9}
(a) $0.1000 \times 15.271 \times 10^{-6}/83.9242 \times 1.3599 \times 10^{-3} = \qty{2.474e-11}{mol}$.
(b) By \eqref{eq:b2:field-lab-methods-2:id}, $N_x = 2.474 \times 10^{-11} \times (732.8 -
0.874)/(0.874 - 0.0568) = \qty{2.216e-8}{mol}$ of $^{86}\mathrm{Sr}$; strontium: $2.216 \times
10^{-8}/0.0986 \times 87.62 = \qty{1.969e-5}{g}$ in \qty{0.0500}{g}, \qty{394}{ppm} (393.9).
(c) A recovery of \qty{99.95}{\percent}: the method is accurate within the rounding of the ratio.
\end{solution}

\begin{solution}{exo:b2:field-lab-methods-2:10}
(a) $R_m = \sqrt{732.8 \times 0.05680} = 6.45$, where $\mathcal M = 1.018$. (b) $\mathcal M = 2.31$
at 0.1, 1.13 at 0.5, 1.038 at 1.664 and 1.16 at 100. (c) Near $R_x$ the mixture hardly differs from
the sample, near $R_s$ from the spike, and the difference that carries the information is lost in
the error of $R_m$: $\mathcal M$ diverges at both ends. Between them the minimum is very flat
($\mathcal M < 1.2$ from about 0.4 to 100), so the analyst only needs the right order of magnitude
of the spike, which a rough concentration, from XRF for instance, provides.
\end{solution}

\begin{solution}{exo:b2:field-lab-methods-2:11}
(a) $\sigma = \sqrt{0.62 \times 0.38/400} = 0.0243$: $62 \pm 2.4$. (b) $(65 - 62)/2.43 = 1.2$
standard deviations: the rock may well be a granodiorite. (c) $\sigma = 0.015$ needs $0.62 \times
0.38/0.015^2 = 1047.1$, so 1048 feldspar points.
\end{solution}

\begin{solution}{exo:b2:field-lab-methods-2:12}
(a) Each point misses with probability $1 - p$, independently: $(1-p)^n$. (b) $0.99^n \le 0.05$
gives $n \ge \ln 0.05/\ln 0.99 = 298.1$: 299 points. (c) $n \ge \ln 0.05/\ln 0.9995 = 5990$: a grid
count is the wrong tool; a rare accessory mineral is searched for by scanning the whole section.
\end{solution}

\begin{solution}{pb:b2:field-lab-methods-2:1}
\noindent\textbf{Part I.}
\textbf{1.} Quartz $372/1200 = \qty{31.0}{\percent}$, alkali feldspar 29.0, plagioclase 27.0, biotite
9.0, muscovite 3.0, opaque oxides 1.0.
\textbf{2.} $\sqrt{0.31 \times 0.69/1200} = 0.0134$; muscovite: $\sqrt{0.97/(0.03 \times 1200)} =
\qty{16.4}{\percent}$.
\textbf{3.} $M = 9 + 3 + 1 = \qty{13.0}{\percent}$; $M' = 13 - 3 = \qty{10.0}{\percent}$ without the
muscovite: leucocratic.
\textbf{4.} $Q + A + P = 1044$ points: Q $= 35.6$, A $= 33.3$, P $= 31.0$.
\textbf{5.} $324/672 = \qty{48.2}{\percent}$, with Q between 20 and 60: a granite, and between 35 and
65 a monzogranite.
\textbf{6.} $\sqrt{0.482 \times 0.518/672} = 0.0193$: the boundaries at 35 and 65 lie 6.9 and 8.7
standard deviations away. The name is secure.
\textbf{7.} $0.31 \times 0.69/0.01^2 = 2139$ points.

\noindent\textbf{Part II.}
\textbf{8.} $0.0051/0.8000 = \qty{0.64}{\percent}$.
\textbf{9.} $0.7949/8.7949 = 0.0904$, and $72 \times 0.0904 = \qty{6.51}{\percent}$ \ce{SiO2} in the
glass: the fusion dilutes the sample elevenfold.
\textbf{10.} $\bar I = 24.91$, $\bar c = 62.283$, $\sum(I - \bar I)^2 = 31.918$, $\sum(I - \bar I)(c -
\bar c) = 80.009$: $a = 2.5067$ per kcps, $b = 62.283 - 2.5067 \times 24.91 = -0.16$.
\textbf{11.} $2.5067 \times 28.63 - 0.16 = \qty{71.61}{\percent}$.
\textbf{12.} BCR-2 $-0.04$, AGV-2 $+0.05$, RGM-1 $-0.01$: all smaller than the standard deviations
of the method (0.1 to 0.2), so the line is adequate over this range.
\textbf{13.} $71.61 + 27.15 + 0.64 = \qty{99.40}{\percent}$, between 97 and 101: not repeated.
\textbf{14.} $71.61 \times 100/(71.61 + 27.15) = \qty{72.51}{\percent}$.

\noindent\textbf{Part III.}
\textbf{15.} $0.0986/0.8258 = 0.1194$: the normalising value is the natural ratio.
\textbf{16.} $\beta = 0.007313/(-0.022972) = -0.3183$.
\textbf{17.} $0.71655 \times (86.908878/85.909261)^{-0.3183} = 0.71392$.
\textbf{18.} The measured 86/88 is $0.11853/0.1194 - 1 = \qty{-0.73}{\percent}$ from the true value
over two mass units: about \qty{0.37}{\percent} per mass unit, in favour of the heavy isotopes.
\textbf{19.} The spike's $^{86}\mathrm{Sr}/^{88}\mathrm{Sr}$ is 0.7661: the mixture's 86/88 is no
longer the natural 0.1194, and fixing it at 0.1194 would treat the spike as an instrumental bias.

\noindent\textbf{Part IV.}
\textbf{20.} Sum of the ratios to 88: $561.4 + 0.7661 + 0.1819 + 1 = 563.348$; $561.4/563.348 =
\qty{99.65}{\percent}$ and $0.7661/563.348 = \qty{0.13599}{\percent}$.
\textbf{21.} $0.2500 \times 15.271 \times 10^{-6}/83.9242 = \qty{4.549e-8}{mol}$ of strontium, and
$\times 1.3599 \times 10^{-3} = \qty{6.186e-11}{mol}$ of $^{86}\mathrm{Sr}$.
\textbf{22.} $561.4/0.7661 = 732.8$ and $0.0056/0.0986 = 0.05680$.
\textbf{23.} $N_x = 6.186 \times 10^{-11} \times (732.8 - 1.6812)/(1.6812 - 0.0568) =
\qty{2.784e-8}{mol}$.
\textbf{24.} $\mathcal M = 1.6812 \times 732.74/(731.12 \times 1.6244) = 1.04$; the ratio adds
$1.04 \times \qty{0.01}{\percent}$, the spike $0.031/15.271 = \qty{0.20}{\percent}$; in quadrature
\qty{0.20}{\percent}, dominated by the calibration of the spike.
\textbf{25.} Strontium: $2.784 \times 10^{-8}/0.0986 = \qty{2.824e-7}{mol}$, that is $2.824 \times
10^{-7} \times 87.62 = \qty{2.474e-5}{g}$ in \qty{0.1000}{g}. The Varden granite contains
\textbf{\qty[text-series-to-math=true]{247}{ppm}} of strontium, $\pm 0.5$.
\end{solution}
